\section{Experimental Setup}
\label{sec:experiments}

\subsection{Hardware Platform}
\label{sec:hardware}

All on-device benchmarks were conducted on a OnePlus~11R (Snapdragon~8+~Gen~1, 8~GB LPDDR5X, UFS~3.1 storage) running Android~14. The device was configured in airplane mode with display brightness fixed at 50\% to minimize thermal and scheduling variability. No other foreground applications were active during benchmarking.

\subsection{Language Model Configuration}
\label{sec:model_config}

The language model used throughout evaluation is Qwen2.5-0.5B-Instruct~\cite{qwen25}, quantized to Q4\_K\_M (4-bit mixed precision) via \texttt{llama.cpp}'s GGUF serialization format. This model was selected as the smallest instruction-tuned model in the Qwen2.5 family capable of producing coherent, context-aware responses on mobile hardware. Greedy sampling was used for all inference trials (temperature~$= 0$, top-$k = 1$) with a maximum generation length of 64 tokens, ensuring deterministic output sequences that minimize generation-phase variance across trials.

\subsection{Benchmark Conditions}
\label{sec:conditions}

The latency evaluation compares three pipeline configurations:

\begin{enumerate}
    \item \textbf{ZERO\_CONTEXT}: The language model receives only the user intent tokens (``\texttt{User: I need water\textbackslash nAssistant:}'') with no environmental context. This establishes the hardware floor for inference latency on the target SoC.

    \item \textbf{RAG\_INLINE}: The full context prompt is concatenated with the user intent and submitted as a single string to \texttt{runInference()}, triggering a complete $O(N)$ token prefill at inference time. This represents the standard retrieval-augmented generation approach. Evaluated at four prompt-length tiers: $N \in \{50, 100, 200, 500\}$ approximate tokens.

    \item \textbf{CAP\_KVC}: A pre-computed KV cache is restored from persistent storage via \texttt{loadKVCache()}, followed by \texttt{resumeInference()} with intent-only tokens. The latency reported includes both the cache restoration time and the subsequent generation time.
\end{enumerate}

The CAP\_KVC condition was evaluated using cache files primed from the base context prompts ($\sim$50 tokens). The N-scaling comparison therefore compares RAG\_INLINE at increasing context lengths against CAP\_KVC at a fixed cache size. This is by design: the CAP\_KVC inference-time cost does not depend on the original context length under the evaluated configuration because the prefill was amortized during background priming. The cache restoration cost is dominated by file I/O (proportional to file size), not by attention computation.

\subsection{Measurement Methodology}
\label{sec:measurement}

Each benchmark condition executes 32 sequential trials: 2~warmup trials (discarded) followed by 30~measured trials. The warmup phase absorbs JIT compilation of Kotlin/ART code paths and initial thermal settling. All measurements are wall-clock latency captured via \texttt{System.nanoTime()} at the Kotlin layer, wrapping the complete JNI call sequence.

\textbf{Important methodological note:} The JNI boundary encapsulates tokenization, prefill, and generation within single native function calls (\texttt{runInference} and \texttt{resumeInference}). True time-to-first-token (TTFT) instrumentation---which requires intercepting the decoder loop between the first and second generated tokens---is not possible without modifying the native inference engine. Accordingly, all latency values reported in this paper represent \textbf{end-to-end inference latency}, encompassing tokenization, prefill, and the complete generation phase (up to 64 tokens). We use the term ``inference latency'' throughout to avoid confusion with TTFT as defined in the LLM serving literature.

For the CAP\_KVC condition, the benchmark instruments two sub-components separately: \texttt{cache\_load\_ms} (the \texttt{loadKVCache} call) and \texttt{inference\_ms} (the \texttt{resumeInference} call), both measured under a single engine lock acquisition to prevent interleaving. The total CAP\_KVC latency is the sum of both components.

Results are emitted as structured CSV lines via Android's \texttt{Log.i()} facility and extracted post-run via \texttt{adb logcat}. The extraction pipeline is deterministic and produces 180 measured data points (6~conditions $\times$ 30~trials).

\subsection{Pre-Benchmark Validation}
\label{sec:validation}

Before executing the benchmark suite, a semantic validation pass confirms that the three pipeline configurations produce qualitatively distinct outputs. A diagnostic prompt (``\texttt{What should I have for breakfast and is my morning medication due?}'') is submitted under each condition, and the generated text is logged for manual inspection. This validates that:

\begin{itemize}
    \item ZERO\_CONTEXT produces a generic response lacking environmental specificity.
    \item RAG\_INLINE produces a response reflecting the \texttt{home\_morning} context (breakfast, medication).
    \item CAP\_KVC produces a similarly context-aware response, confirming that KV cache restoration successfully preserved the contextual attention patterns.
\end{itemize}

The benchmark suite is aborted if the KV cache load fails during validation, preventing the collection of meaningless data.

\subsection{Memory Profiling}
\label{sec:memory_profiling}

Native heap memory usage is measured using Android's \texttt{dumpsys meminfo} command at two states: (1)~after model initialization with zero KV cache states loaded, and (2)~after loading three KV cache states into native sequence slots. The delta represents the marginal memory cost of the predictive caching system.

\subsection{Fusion Model Evaluation}
\label{sec:fusion_eval}

The multimodal fusion model is evaluated through a controlled ablation study comparing two architectures:

\begin{itemize}
    \item \textbf{Cross-Attention}: A single-head cross-attention module ($d = 64$) where the EMG embedding serves as the query and the projected gaze vector serves as key and value.
    \item \textbf{Late Fusion}: Concatenation of the 64-dimensional EMG embedding and the 64-dimensional projected gaze vector, followed by a two-layer MLP classifier.
\end{itemize}

Both architectures are trained on a synthetic dataset of 2000 paired EMG-gaze samples (60\% train, 20\% validation, 20\% test) drawn from a synthetic data generator that pairs random EMG embeddings with corresponding gaze vectors. The pre-specified decision criterion requires that the cross-attention model exceed the late fusion model's macro F1 score by more than 3~percentage points (absolute) to justify deployment of the more complex architecture.

\textbf{Limitation:} The fusion training data is synthetically generated and does not represent real paired EMG-gaze recordings from AAC users. The reported accuracy and F1 scores characterize the model's capacity to learn the fusion mapping under controlled conditions, not its expected clinical accuracy. This is explicitly acknowledged as a limitation of the current evaluation.

\subsection{EMG Classifier Evaluation}
\label{sec:emg_eval}

The CNN-LSTM EMG classifier is trained on a subset of the NinaPro~DB5 dataset~\cite{ninapro} using a leave-one-subject-out (LOSO) protocol: subjects~1--9 for training and subject~10 for validation. Training uses the Adam optimizer with learning rate $10^{-3}$, cosine annealing schedule, class-weighted cross-entropy loss, and gradient clipping at norm~1.0. The model is trained for up to 50~epochs with early stopping based on validation accuracy.

The EMG classifier operates as an embedding extractor for the downstream fusion model. Its standalone classification accuracy on the held-out subject is reported separately from the fusion accuracy to avoid conflation.

\subsection{Drift Detection Ablation}
\label{sec:drift_eval}

The sensitivity of conversation-level drift detection to the sliding window size $k$ is evaluated offline on the DailyDialog dataset~\cite{dailydialog}. Each conversation is annotated with ground-truth topic shift indices. For each $k \in \{1, 2, 3, 4, 5\}$, the system computes the centroid embedding of the most recent $k$ turns using Sentence-BERT (\texttt{all-MiniLM-L6-v2}~\cite{minilm}), and predicts a topic shift when the cosine similarity between the anchor embedding and the centroid falls below a threshold $\theta = 0.15$. Precision, recall, and F1 are computed for each $k$. The MiniLM model is used solely for this offline ablation analysis and is not part of the deployed Android runtime; the on-device drift detector operates on router score deltas rather than semantic embeddings.

\textbf{Limitation:} DailyDialog is a general-purpose dialogue corpus that differs substantially from AAC communication patterns in terms of turn length, vocabulary, and topic variability. The k-ablation results characterize the detector's behavior on available annotated data rather than providing direct evidence of clinical effectiveness.
